\documentclass[11pt]{article}
\usepackage[margin=1in]{geometry}
\usepackage[T1]{fontenc}
\usepackage{amsmath,amssymb,booktabs,tabularx,array}
\usepackage[hidelinks]{hyperref}
\newcolumntype{Y}{>{\raggedright\arraybackslash}X}
\newcommand{\code}[1]{\texttt{\detokenize{#1}}}
\title{Computational Complexity of TrIM: A Six-Step Analysis}
\author{Nina Li}
\date{}
\begin{document}
\maketitle

\noindent This note follows the current Mondrian-based implementation in
\code{estimatedH_CSS/src/}. The main analysis covers one initial EGOP estimate,
followed by downstream regression. Additional iterations and the experimental
lifetime sweep are identified separately. No experiments were run.

\paragraph{Notation.}
All time bounds refer to the current implementation, using conventional dense
arithmetic costs. The symbols below are used throughout.
\begin{center}
\small
\begin{tabularx}{\textwidth}{lY}
\toprule
Symbol & Definition \\
\midrule
$n$ & Number of training samples. \\
$q$ & Number of test samples. \\
$d$ & Number of input coordinates; the full TrIM transform retains $d$ columns. \\
$M$ & Number of Mondrian trees in a forest. \\
$\lambda$ & Maximum lifetime, which determines which sampled split events occur. \\
$m$, $e$ & Tree index and processed split-event index, respectively. \\
$S_m$, $S$ & Splits in tree $m$ and total forest splits, $S=\sum_{m=1}^{M}S_m$. \\
$L_m$, $C$ & Active leaves in tree $m$, $L_m=1+S_m$; total node-column count $C=M+2S$. \\
$n_e$ & Number of training samples in the node processed at split $e$. \\
$B$ & Training-sample participation across splits: $B=\sum_{e=1}^{S}n_e$. \\
$t$, $B_{f,t}$ & Number of samples in a prediction batch, and their total participation across replayed splits in forest $f$. In particular, $B_{1,q}$ refers to the downstream forest's test batch. \\
$F$, $P(t)$ & Time for one forest fit and one batch prediction, expanded below. \\
$f$ & Forest label: $f=0$ for the initial forest and $f=1$ for the downstream forest. A subscript $f$ attaches split counts and costs to that forest. \\
$N$, $J$, $W$ & Total samples $N=n+q$; checkpoint count $J=\lceil S_1/50\rceil$ for $S_1>0$; combined downstream lifetime-sweep time $W$. \\
\bottomrule
\end{tabularx}
\end{center}

\section{Train the Initial Forest}

\code{MondrianForestRegressor.fit} calls \code{train} in
\code{Mondrian_RF/utils.py}. The operations below describe one forest fit on
the original training inputs and responses; each time bound covers all trees
unless stated otherwise.

\subsection*{1.1 Initialize data and memberships}
The code copies the inputs, centers the responses, and computes the initial
coordinate bounds shared by all trees. It also records that every sample
initially belongs to each tree's root.
\[
\text{Time: }O(nd+nM).
\]
Here $n$ is the number of training samples, $d$ is the number of coordinates,
and $M$ is the number of trees. Scanning and copying inputs costs $O(nd)$;
initializing one membership per sample per tree costs $O(nM)$.

\subsection*{1.2 Sample and schedule cuts}
The code samples root cuts and two child candidates per processed split,
using coordinate ranges to choose their times and locations. A heap processes
accepted events in time order, before the maximum lifetime is reached.
\[
\text{Time: }O\bigl(Md+Sd+(M+S)\log(1+M+S)\bigr).
\]
Here $M$ is the number of trees, $d$ is the number of coordinates, and $S$ is
the total processed split count across the forest. Root sampling costs $O(Md)$,
child sampling costs $O(Sd)$, and the remaining term bounds heap operations.
Root sampling is performed even for trees whose first candidate is rejected.

\subsection*{1.3 Split nodes and update records}
Each split partitions its samples, updates memberships, and recomputes child
bounds. The code also removes the parent from its tree's active-leaf list
and appends the children; Python list removal has linear cost.
\[
\text{Time: }O\!\left(dB+S+\sum_{m=1}^{M}S_m^2\right),
\qquad B=\sum_{e=1}^{S}n_e.
\]
Here $d$ is the coordinate count, $M$ is the tree count, $m$ indexes trees,
$S_m$ is the number of splits in tree $m$, and $S=\sum_m S_m$.
The index $e$ identifies a processed split, $n_e$ is its node's training-sample
count, and $B$ is total sample participation across splits. The $dB$ term
covers coordinate scans; the remaining terms bound active-list maintenance.

\emph{Tree shape and lifetime.} With $n$ training samples and $M$ trees,
balanced growth gives $B=O(Mn\log n)$; an unbalanced-tree upper bound is
$O(Mn^2)$. Balanced growth is not guaranteed. Assuming nonempty children,
$S\leq M(n-1)$ and $B\geq2S$, so the $Sd$ sampling term in Step 1.2 is absorbed
by $dB$ in the total bound. Maximum lifetime $\lambda$ changes which splits
occur, and hence $S$ and $B$, rather than introducing a universal linear
runtime factor in $\lambda$.

\subsection*{1.4 Extract memberships and compute leaf weights}
After growth, the code extracts each tree's active columns from the shared
CSR matrix. The optimized \code{_leaf_weights} then computes the diagonal and
right-hand side and divides elementwise, without a dense linear solve.
\[
\begin{aligned}
\text{Extraction per tree (conditional): }&
O\bigl(nM+C+L_m\log(1+L_m)\bigr),\\
\text{Weight calculation across all trees: }&O(nM+M+S).
\end{aligned}
\]
Here $n$ is the training-sample count, $M$ is the tree count, $S$ is the total
split count, $C=M+2S$ is the matrix's node-column count, and $L_m=1+S_m$ is
the active-leaf count in tree $m$, which has $S_m$ splits.
Across all trees, extraction contributes $O(nM^2+MC)$ plus sorting.
Since $\sum_m L_m=M+S$, sorting is bounded by $(M+S)\log(1+M+S)$ and can
share the total bound used for heap operations; the heap does not perform it.

\subsection*{1.5 Total time and working memory}
Adding these operations gives the following conditional fit-time upper bound:
\begin{equation}
\boxed{\begin{aligned}
F=O\!\bigl(&nd+nM^2+MC+Md+dB\\
&+(M+S)\log(1+M+S)+\sum_{m=1}^{M}S_m^2\bigr).
\end{aligned}}
\end{equation}
Here $F$ is one forest's fit time; $n$ is the training-sample count; $d$ is the
coordinate count; $M$ is the tree count; $S_m$ is the split count in tree $m$;
$S=\sum_m S_m$; $C=M+2S$; and $B=\sum_e n_e$, where $n_e$ counts the training
samples in the node processed at split $e$.

The working-memory upper bound is
\[
O(nd+nM+M+S+B),
\]
where $n$ counts training samples, $d$ coordinates, $M$ trees, $S$ total splits,
and $B$ total training-sample participation. This includes input arrays,
memberships, node/event records, and retained parent/child sample-index arrays;
temporary construction arrays are released after \code{train} returns.

For the initial forest, later sections write $F_0$, $S_0$, $C_0$, $B_0$, and
$S_{0,m}$ for these quantities. More generally, $f$ labels a forest,
$S_{f,m}$ is its tree $m$'s split count, and $C_f=M+2S_f$.

\section{Estimate the EGOP Matrix}

\subsection*{2.1 Create perturbed inputs and estimate gradients}
For every coordinate, \code{TrIM._get_importance} copies the training inputs
and applies positive and negative perturbations. It queries the initial forest
on both batches and divides their prediction difference by the step size.
\[
\text{Time: }O\bigl(nd^2+dP_0(n)\bigr).
\]
Here $n$ is the training-sample count, $d$ is the coordinate count, and $P_0(n)$
is a common upper bound for one initial-forest prediction on an $n$-sample
perturbed batch. Copying full arrays across coordinates costs $O(nd^2)$;
there are $2d$ prediction calls, whose cost is expanded in Section 5.

\emph{Different routes.} Perturbations can change sample routes. The bound
$P_0(n)$ uses the maximum routing participation across these batches, so it
does not assume identical or similar routes.

\subsection*{2.2 Form the EGOP matrix}
The code stacks the estimated gradients and computes their outer-product
average using one dense matrix multiplication.
\[
H=GG^\top/n,\qquad \text{Time: }O(nd^2).
\]
Here $n$ is the training-sample count, $d$ is the coordinate count,
$G\in\mathbb{R}^{d\times n}$ contains one gradient per column, and
$H\in\mathbb{R}^{d\times d}$ is the estimated EGOP matrix. The product computes
$d^2$ entries, each summing over $n$ samples.

\subsection*{2.3 Total time and additional iterations}
Adding perturbation, prediction, and matrix formation gives the first-round
EGOP time, excluding the initial forest fit counted in Section 1:
\[
\boxed{T_2=O\bigl(nd^2+dP_0(n)\bigr).}
\]
Here $T_2$ is the first EGOP round's time, $n$ is the training-sample count,
$d$ is the coordinate count, and $P_0(n)$ bounds one initial-forest prediction
on any of the perturbed training batches.

In later rounds, every perturbed batch is independently transformed before
prediction. The resulting bound per EGOP estimate is
\[
O\bigl(nd^3+dP(n)\bigr),
\]
where $n$ counts training samples, $d$ coordinates, and $P(n)$ bounds one
prediction by that round's fitted forest. The $2d$ dense transformations each
cost $O(nd^2)$; the preceding forest refit is a separate cost.

Gradient storage, perturbed batches, and the EGOP matrix require
$O(nd+d^2)$ additional space, where $n$ counts training samples and $d$
coordinates. Forest-prediction working memory is additional; calls run
sequentially, so peak memory is not multiplied by the number of calls.

\section{Transform the Inputs}

\subsection*{3.1 Normalize the EGOP matrix}
\code{TrIM.transform} divides the EGOP matrix by the mean of its column norms.
This normalization is recomputed whenever the method is called.
\[
A=H/a(H),\qquad a(H)=\frac{1}{d}\sum_{j=1}^{d}\lVert H_{:,j}\rVert_2,
\qquad \text{Time: }O(d^2).
\]
Here $H$ is the $d\times d$ EGOP matrix, $d$ is the coordinate count, $j$
indexes columns, $a(H)$ is the normalization scalar, and $A$ is the normalized
matrix. The transformation assumes $a(H)\ne0$.

\subsection*{3.2 Transform training and test inputs}
The code multiplies each input batch by the normalized matrix. Responses
remain unchanged, and the full transformation retains the input dimension.
\[
Z=XA,\qquad \boxed{T_3=O\bigl((n+q)d^2\bigr).}
\]
Here $X$ is an input batch, $A$ is the $d\times d$ normalized EGOP matrix,
$Z$ is the transformed batch, $n$ is the training-sample count, $q$ is the
test-sample count, and $d$ is the coordinate count. $T_3$ includes both batch
multiplications and their normalizations.

Output arrays and the transformation require $O((n+q)d+d^2)$ space, where
$n$ and $q$ count training and test samples and $d$ counts coordinates.

\section{Train the Downstream Forest}

\subsection*{4.1 Fit at a fixed lifetime}
The code trains a new forest on transformed training inputs and the original
responses. Its construction operations are the same as in Section 1, but
transformation can change the resulting split counts and sample participation.
\[
\begin{aligned}
\text{Time: }F_1=O\bigl(&nd+nM^2+MC_1+Md+dB_1\\
&+(M+S_1)\log(1+M+S_1)+\sum_{m=1}^{M}S_{1,m}^2\bigr).
\end{aligned}
\]
Here $F_1$ is downstream fit time, $n$ counts training samples, $d$ coordinates,
$M$ trees, and $m$ indexes trees. $S_{1,m}$ counts splits in downstream tree
$m$, $S_1=\sum_m S_{1,m}$, $C_1=M+2S_1$ counts node columns, and $B_1$ is total
training-sample participation across downstream splits. Sparse extraction
uses the conditional bound described in Section 1.

Working memory is $O(nd+nM+M+S_1+B_1)$: $n$ counts training samples, $d$
coordinates, $M$ trees, $S_1$ downstream splits, and $B_1$ training-sample
participation. Prediction and error evaluation follow in Sections 5--6.

\subsection*{4.2 Alternative: grow a forest through lifetime checkpoints}
The experiments instead use \code{evaluate_all_lifetimes}, which interleaves
training, prediction, and MSE calculation. It grows one forest rather than
independently fitting a forest at every checkpoint.
\[
\begin{aligned}
\text{Growth time: }O\bigl(&Nd+NM+Md+dB_1+B_{1,q}\\
&+(M+S_1)\log(1+M+S_1)+\sum_{m=1}^{M}S_{1,m}^2\bigr).
\end{aligned}
\]
Here $N=n+q$ combines $n$ training and $q$ test samples; $d$ is the coordinate
count and $M$ the tree count. $S_{1,m}$ counts splits in downstream tree $m$,
$S_1=\sum_m S_{1,m}$, and $B_1$ and $B_{1,q}$ are training- and test-sample
participation across splits. Child bounds use only training samples; test
samples are routed without scanning all their coordinates at each split.

\subsection*{4.3 Checkpoint work and total sweep cost}
At each checkpoint, the code extracts the current tree's memberships,
recomputes its weights and predictions, averages stored tree predictions,
and computes MSE.
\[
\text{Checkpoint work: }O\bigl(J[NM+C_1+n\log(1+n)]\bigr).
\]
Here $J=\lceil S_1/50\rceil$ is the checkpoint count for $S_1>0$ processed
splits, $N=n+q$ combines training and test counts, $M$ is the tree count, and
$C_1=M+2S_1$ is the node-column count. Checkpoints occur after splits
$1,51,101,\ldots$. The $NM$ term covers membership scans and prediction
averaging; $n\log(1+n)$ bounds sorting selected leaf indices, assuming each
leaf contains a training sample.

Adding growth and checkpoint costs gives the conditional total sweep time:
\[
\boxed{\begin{aligned}
W=O\bigl(&Nd+NM+Md+dB_1+B_{1,q}\\
&+(M+S_1)\log(1+M+S_1)+\sum_{m=1}^{M}S_{1,m}^2\\
&+J[NM+C_1+n\log(1+n)]\bigr).
\end{aligned}}
\]
Here $W$ combines Steps 4--6; $N=n+q$ counts all samples, $d$ coordinates,
$M$ trees, $S_{1,m}$ splits in tree $m$, $S_1=\sum_m S_{1,m}$ total splits,
$C_1=M+2S_1$ node columns, $B_1$ training participation, $B_{1,q}$ test
participation, and $J=\lceil S_1/50\rceil$ checkpoints.
Working memory is $O(Nd+NM+M+S_1+B_1+B_{1,q})$ with the same quantities.
Do not add separate fixed-lifetime fitting and prediction costs when using $W$.

\emph{Checkpoint behavior.} Only the tree associated with the current event
is refreshed; other prediction columns retain earlier values and initially
contain zeros. There is no forced terminal checkpoint, so this is not a
synchronized evaluation of every tree. The bound describes successful runs
with splits; the zero-split case reaches an unassigned prediction variable.

\section{Predict on the Test Inputs}

\subsection*{5.1 Initialize prediction data and replay splits}
For fixed-lifetime prediction, \code{evaluate} initializes memberships and
replays the stored split history on the transformed inputs. It uses
\code{history.pop(0)}, which repeatedly shifts the remaining list entries.
\[
\text{Time: }O\bigl(n+td+tM+B_{f,t}+S_f^2\bigr).
\]
Here $n$ counts training responses whose mean is recomputed, $t$ is the
prediction-batch size, $d$ the coordinate count, and $M$ the tree count.
The label $f$ identifies the queried forest: $0$ initial, $1$ downstream.
$S_f$ counts its splits, and $B_{f,t}$ counts the batch's sample participation
across replayed splits. List-front removal alone costs $\Theta(S_f^2)$;
this upper bound also covers active-leaf list updates.

\subsection*{5.2 Extract memberships and average tree predictions}
The code extracts each tree's active columns, applies the stored leaf weights,
and averages tree predictions. Input transformation is already counted in
Section 3.
\[
\text{Conditional time: }O\bigl(tM^2+MC_f+S_f^2\bigr).
\]
Here $t$ is the batch size, $M$ the tree count, $S_f$ the queried forest's
split count, and $C_f=M+2S_f$ its node-column count.
The extraction bound includes sorting: summed over trees, its cost is
$O((M+S_f)^2)=O(MC_f+S_f^2)$. Applying weights and averaging costs $O(tM)$,
which is covered by the displayed bound.

\subsection*{5.3 Total prediction time and memory}
Combining initialization, replay, extraction, and averaging gives
\[
\boxed{P_f(t)=O\bigl(n+td+tM^2+MC_f+B_{f,t}+S_f^2\bigr).}
\]
Here $P_f(t)$ is batch-prediction time for forest $f$, $n$ is the training
count, $t$ the batch size, $d$ the coordinate count, $M$ the tree count,
$C_f=M+2S_f$ the node-column count, $S_f$ the split count, and $B_{f,t}$ total
batch-sample participation. Use $f=1,t=q$ for $q$ test samples; use $f=0,t=n$
for initial EGOP queries, taking maximum participation across perturbations
for their common upper bound.

Working memory is $O(td+tM+M+S_f+B_{f,t})$, where $t$ counts batch samples,
$d$ coordinates, $M$ trees, $S_f$ splits, and $B_{f,t}$ routing participation.
The stored model is additional. Sweep prediction is already included in $W$,
the combined lifetime-sweep time from Section 4.

\section{Evaluate Prediction Error}

\subsection*{6.1 Calculate mean squared error}
The code compares predictions with the true test responses and averages the
squared differences. Test responses are used for scoring, not for fitting
leaf weights.
\[
\mathrm{MSE}=\frac{1}{q}\sum_{i=1}^{q}(\widehat y_i-y_i)^2,
\qquad \text{Time: }O(q).
\]
Here $q$ is the test-sample count, $i$ indexes test samples, $y_i$ is a true
response, and $\widehat y_i$ is its prediction. Vectorized evaluation uses at
most $O(q)$ temporary space.

\subsection*{6.2 Record and summarize checkpoint errors}
For lifetime evaluation, the code records errors at checkpoints; comparison
helpers can then report the minimum recorded test MSE.
\[
\text{MSE calculations: }O(Jq),\qquad
\text{Minimum search: }O(J).
\]
Here $J$ is the number of checkpoints and $q$ the test-sample count. Recorded
times and errors occupy $O(J)$ space. Checkpoint MSE work is already included
in the sweep cost; reporting the minimum is not an independent validation
and final-refit procedure.

\subsection*{6.3 Total cost of the six-step pipeline}
For one EGOP round followed by fixed-lifetime downstream regression, add the
six steps without counting any fit or transformation twice:
\[
\boxed{\begin{aligned}
T_{\mathrm{total}}=F_0+F_1
+O\bigl(&dP_0(n)+(n+q)d^2\\
&+P_1(q)+q\bigr).
\end{aligned}}
\]
Here $T_{\mathrm{total}}$ is total time, $F_0$ and $F_1$ are initial and
downstream forest-fit times, $P_0(n)$ bounds an initial-forest prediction on
any perturbed training batch, and $P_1(q)$ is downstream prediction time.
The symbols $n$, $q$, and $d$ count training samples, test samples, and
coordinates, respectively.

For the experimental lifetime sweep, use its combined cost in place of
separate downstream fitting, prediction, and MSE costs:
\[
\boxed{T_{\mathrm{total,sweep}}
=F_0+O\bigl(dP_0(n)+(n+q)d^2+W\bigr).}
\]
Here $T_{\mathrm{total,sweep}}$ is total time with lifetime evaluation,
$F_0$ is initial forest-fit time, $P_0(n)$ is the common initial prediction
bound, $n$ counts training samples, $q$ test samples, $d$ coordinates, and
$W$ combines downstream growth, checkpoint prediction, and error evaluation.
These totals exclude data generation, oracle and CSS comparisons, repeated
trials, and additional EGOP rounds. Actual runtime shares require profiling.

\paragraph{Source references.}
In \code{Mondrian_RF/Mondrian_forest.py}, see lines 32--39 for forest wrappers,
83--108 for EGOP estimation, and 110--136 for TrIM fitting and transformation.
In \code{Mondrian_RF/utils.py}, see lines 8--33 for leaf weights, 36--130 for
training, 132--184 for prediction, and 209--325 for the lifetime sweep.

\end{document}
